\chapter{MongoDB Aggregation}

\section{The Pipeline Concept}

An aggregation pipeline processes documents through a sequence of stages,
each one transforming the output of the previous stage. Unlike
\texttt{.find()}, which returns matching documents as-is, aggregation can
reshape, group, and compute across documents --- exactly what's needed for
dashboard statistics.

\begin{diagrambox}[Aggregation Pipeline]
\begin{verbatim}
Task collection
      |
      v
  $match   --> filter documents (like a WHERE clause)
      |
      v
  $group   --> bucket documents and compute per-bucket values
      |
      v
  $sort    --> order the grouped results
      |
      v
  $project --> reshape the output fields
      |
      v
  Final result array
\end{verbatim}
\end{diagrambox}

\section{Dashboard Stats Endpoint}

\texttt{GET /api/tasks/stats} needs: total task count, completed count,
pending count, and a breakdown of tasks per user. This maps to two separate
\texttt{\$group} stages in one pipeline.

\begin{lstlisting}[language=JavaScript]
// backend/src/controllers/task.controller.js
const Task = require("../models/Task");

const getTaskStats = async (req, res, next) => {
  try {
    const statusCounts = await Task.aggregate([
      // Stage 1: only this user's tasks (skip if stats are app-wide)
      { $match: { owner: req.user.id } },

      // Stage 2: group by status, counting documents in each group
      {
        $group: {
          _id: "$status",
          count: { $sum: 1 },
        },
      },

      // Stage 3: alphabetical by status for a stable response order
      { $sort: { _id: 1 } },
    ]);

    // Reshape [{ _id: "completed", count: 5 }, ...] into a flat object
    const summary = { total: 0, completed: 0, pending: 0, "in-progress": 0 };
    statusCounts.forEach(({ _id, count }) => {
      summary[_id] = count;
      summary.total += count;
    });

    const perUser = await Task.aggregate([
      {
        $group: {
          _id: "$owner",
          taskCount: { $sum: 1 },
        },
      },
      // Join back to the users collection to attach a readable name
      {
        $lookup: {
          from: "users",
          localField: "_id",
          foreignField: "_id",
          as: "user",
        },
      },
      { $unwind: "$user" },
      {
        $project: {
          _id: 0,
          userId: "$_id",
          name: "$user.name",
          taskCount: 1,
        },
      },
      { $sort: { taskCount: -1 } },
    ]);

    res.status(200).json({ summary, perUser });
  } catch (error) {
    next(error);
  }
};

module.exports = { getTaskStats };
\end{lstlisting}

\begin{notebox}[NOTE]
\texttt{\$lookup} performs a join against another collection --- here,
pulling in each owner's \texttt{name} from the \texttt{users} collection so
the response doesn't just show raw \texttt{ObjectId}s. It plays the same
role that \texttt{.populate()} plays in a normal Mongoose query, but inside
an aggregation pipeline \texttt{.populate()} itself isn't available, so
\texttt{\$lookup} is the aggregation-native equivalent.
\end{notebox}

\subsection{Example Output}

\begin{lstlisting}[language=JavaScript]
{
  "summary": {
    "total": 45,
    "completed": 20,
    "pending": 15,
    "in-progress": 10
  },
  "perUser": [
    { "userId": "64f...a1", "name": "Asha Patel", "taskCount": 18 },
    { "userId": "64f...b2", "name": "Ravi Kumar", "taskCount": 15 },
    { "userId": "64f...c3", "name": "Meera Nair", "taskCount": 12 }
  ]
}
\end{lstlisting}

\section{Route Wiring}

\begin{lstlisting}[language=JavaScript]
// backend/src/routes/task.routes.js
router.get("/stats", protect, getTaskStats);
\end{lstlisting}

\begin{tipbox}[TIP]
Aggregation stages run in the order you list them, and \texttt{\$match}
stages placed early narrow down the working set before the more expensive
\texttt{\$group}/\texttt{\$lookup} stages run --- put filters as close to the
top of the pipeline as possible.
\end{tipbox}
